\section{The exact rank of the denominator obstruction}
\label{jnew:sec:rank}

The single-value criterion also yields a complete finite-dimensional
calculation for the denominator part of the companion. Let $q$ be odd,
$G=G_q$, and $H=\langle2\rangle\subseteq G$. For a finite abelian
group $U$, write
\[
 d(U)=\frac{|U|-|U[2]|}{2}.
\]
The integer $d(U)$ is the dimension of the subspace of
$\mathbb Q[U]$ that changes sign under inversion.
For the boundary case $q=1$, use $G_1=H_1=\{1\}$ and $\rho(1)=0$;
the symbol families are zero, so the assertions below are vacuous there.

\begin{theorem}[Rank and all coefficient relations]
\label{jnew:thm:denominator-rank}
The two subspaces
\[
 \operatorname{span}_{\mathbb Q}\{D_q(a):a\in G\},
 \qquad
 \operatorname{span}_{\mathbb Q}\{T_q(a):a\in G\}
\]
are equal. Their common dimension is
\begin{equation}\label{jnew:eq:denominator-rank}
 \boxed{\displaystyle\rho(q)=d(G)-d(G/H)
 =\frac{|G|-|G[2]|-|G/H|+|(G/H)[2]|}{2}.}
\end{equation}
For rational coefficients $c_a$, the exact relation test for the
second-difference family is
\begin{equation}\label{jnew:eq:T-relations}
 \sum_{a\in G}c_aT_q(a)=0
 \quad\Longleftrightarrow\quad
 c_a-c_{a^{-1}}\text{ is constant on each coset of }H.
\end{equation}
For the first-difference family the exact relation test is
\begin{equation}\label{jnew:eq:D-relations}
 \sum_{a\in G}c_aD_q(a)=0
 \quad\Longleftrightarrow\quad
 c_a-c_{2a}=c_{a^{-1}}-c_{2a^{-1}}\quad(a\in G).
\end{equation}
\end{theorem}
\begin{proof}
The existing all-conductor kernel theorem identifies the symbol span
with the inversion-odd subspace $V^-$ of $\mathbb Q[G]$, by
$\beta_q(a)\mapsto[a]-[a^{-1}]$. Work first over $\mathbb R$, with
its usual inner product, and write $P_g[a]=[ga]$, where $g$ is the
class of $2$. The cyclic Laplacian
$L_g=P_g+P_{g^{-1}}-2I$ is self-adjoint, preserves $V^-$, and has
kernel the functions constant on $H$-cosets. Thus the span of the
$T_q(a)$ is the orthogonal complement in $V^-$ of
$(V^-)^H$. The latter has dimension $d(G/H)$, proving
\eqref{jnew:eq:denominator-rank} for $T$ and the coefficient test
\eqref{jnew:eq:T-relations}.

An inversion-odd vector $w$ is orthogonal to every
$D_q(a)=([a]-[a^{-1}])-([a/g]-[ga^{-1}])$ precisely when
\[
 2(w_a-w_{a/g})=0\qquad(a\in G).
\]
This is exactly the condition that $w$ be constant on $H$-cosets.
Hence the two spans have the same orthogonal complement and are
equal. All matrices have rational entries, so these dimensions and
equalities descend to $\mathbb Q$. Finally, the coefficient of
$[a]$ in the inversion-odd image of a $D$-combination is
$c_a-c_{ga}-c_{a^{-1}}+c_{ga^{-1}}$, proving
\eqref{jnew:eq:D-relations}.
\end{proof}

For example, at $q=31$, $G$ is cyclic of order $15$ and the class
of $2$ has order $5$. Therefore
\[
 d(G)=7,\qquad d(G/H)=d(C_3)=1,\qquad\rho(31)=6.
\]
The loss of one dimension is an exact distribution relation, not a
numerically small singular value. The relations include the evident
telescoping identities
\[
 \sum_{h\in H}D_q(ah)=0,\qquad
 \sum_{h\in H}T_q(ah)=0,
\]
and the theorem describes all additional consequences of inversion.

\begin{corollary}[A certified lower bound for formal companion values]
\label{jnew:cor:value-rank}
For each odd $q$, there exist $\rho(q)$ positive even numerators
$e_j$, coprime to $q$, such that the formal elements
$\eta_{e_j/q}$ are linearly independent modulo rational
dilogarithms.
\end{corollary}
\begin{proof}
Choose residues giving a basis among $T_q(a)$, and choose even
positive representatives $e_j$ of those residues. In a common
cyclotomic field, apply the exterior square of the normalized norm
to $\mathbb Q(\zeta_q)$. It kills every numerator-conductor block
and keeps the corresponding $T_q(e_j)$. If a combination of the
$\eta_{e_j/q}$ belonged to the rational pre-Bloch group, this
projection would be a rational wedge. The normalized norm to
$\mathbb Q$ kills the symbol combination and fixes that wedge, so
the wedge is zero. The chosen $T$ symbols are independent, proving
the assertion.
\end{proof}
